\section{Recursion}

\textbf{Recursion} is a programming technique where a function calls itself to solve a problem by breaking it into smaller subproblems.  
Every recursive function must have a \textbf{base case} to stop the recursion.

\subsection{Memory Model}
Each function call (recursive or not) creates a new \textbf{stack frame} on the call stack, which stores:
\begin{itemize}
  \item local variables
  \item parameters
  \item return address
\end{itemize}
Recursive calls create multiple stack frames, which are removed as the function returns. Recursive depth determines stack memory usage.

\subsection{Example: Factorial}

Computing factorials with a for loop:
\begin{verbatim}
int factorial(int n) {
    int result = 1;
    for (int i = 1; i <= n; i++) {
        result *= i;
    }
    return result;
}
\end{verbatim}
Computing factorials with recursion:
\begin{verbatim}
int factorial(int n) {
    if (n <= 1) {
        return 1;   // base case
    }
    return n * factorial(n - 1);
}
\end{verbatim}

\subsection{Example: Fibonacci}
Computing fibonacci with recursion:
\begin{verbatim}
int fib(int x) {
    if (x == 1) {
        return 1;
    }
    if (x <= 0) {
        return 0;
    }
    return fib(x - 1) + fib(x - 2);
}
\end{verbatim}